\documentclass[12pt,a4paper]{report}

% ============================================================
% 1. FONT VÀ TIẾNG VIỆT - NÊN BIÊN DỊCH BẰNG XeLaTeX
% ============================================================
\usepackage{fontspec}
\usepackage{polyglossia}
\setmainlanguage{vietnamese}

% Ưu tiên Times New Roman nếu máy có.
% Nếu không có Times New Roman thì dùng TeX Gyre Termes (font gần giống Times).
% Nếu cũng không có thì dùng Latin Modern Roman.
\defaultfontfeatures{Mapping=tex-text}
\IfFontExistsTF{Times New Roman}{%
  \setmainfont{Times New Roman}%
}{%
  \IfFontExistsTF{TeX Gyre Termes}{%
    \setmainfont{TeX Gyre Termes}%
  }{%
    \setmainfont{Latin Modern Roman}%
  }%
}

% ============================================================
% 2. ĐỊNH DẠNG TRANG, GIÃN DÒNG, THỤT ĐẦU DÒNG
% ============================================================
\usepackage[left=3cm,right=2cm,top=2.5cm,bottom=2.5cm]{geometry}
\usepackage{setspace}
\usepackage{indentfirst}

\onehalfspacing
\setlength{\parindent}{1cm}
\setlength{\parskip}{0pt}

% ============================================================
% 3. CÁC GÓI HỖ TRỢ HÌNH ẢNH, BẢNG, CHÚ THÍCH, LIÊN KẾT
% ============================================================
% [demo] giúp biên dịch được ngay cả khi chưa có file ảnh.
% Khi đã có ảnh thật, hãy bỏ [demo].
\usepackage[demo]{graphicx}
\graphicspath{{images/}}

\usepackage{caption}
\usepackage{booktabs}
\usepackage{array}
\usepackage{tabularx}
\usepackage{longtable}
\usepackage{multirow}
\usepackage{amsmath}
\usepackage{float}
\usepackage{tocloft}
\usepackage{titlesec}
\usepackage[hidelinks]{hyperref}

% Caption của bảng nằm phía trên bảng.
% Caption của hình nằm phía dưới hình.
\captionsetup[table]{
  position=above,
  skip=6pt,
  justification=centering,
  labelfont=bf,
  font=small
}

\captionsetup[figure]{
  position=below,
  skip=6pt,
  justification=centering,
  labelfont=bf,
  font=small
}

% ============================================================
% 4. ĐÁNH SỐ CÔNG THỨC, HÌNH, BẢNG THEO CHƯƠNG
% Ví dụ: Hình 2.1, Bảng 3.2, Công thức (2.1)
% ============================================================
\numberwithin{equation}{chapter}
\numberwithin{figure}{chapter}
\numberwithin{table}{chapter}

% ============================================================
% 5. TÊN CÁC PHẦN BẰNG TIẾNG VIỆT
% ============================================================
\AtBeginDocument{%
  \renewcommand{\chaptername}{Chương}
  \renewcommand{\contentsname}{MỤC LỤC}
  \renewcommand{\listfigurename}{DANH MỤC HÌNH VẼ}
  \renewcommand{\listtablename}{DANH MỤC BẢNG}
  \renewcommand{\bibname}{TÀI LIỆU THAM KHẢO}
  \renewcommand{\figurename}{Hình}
  \renewcommand{\tablename}{Bảng}
}

% ============================================================
% 6. ĐỊNH DẠNG CHƯƠNG, MỤC, MỤC LỤC
% ============================================================
\setcounter{secnumdepth}{3}
\setcounter{tocdepth}{3}

% Chương hiển thị dạng:
% Chương 1. TỔNG QUAN ĐỀ TÀI
\titleformat{\chapter}[block]
  {\normalfont\Large\bfseries\centering}
  {\chaptertitlename\ \thechapter.}{0.5em}{\MakeUppercase}
\titlespacing*{\chapter}{0pt}{0pt}{20pt}

% Section, subsection, subsubsection
\titleformat{\section}
  {\normalfont\bfseries}
  {\thesection}{0.6em}{}

\titleformat{\subsection}
  {\normalfont\bfseries}
  {\thesubsection}{0.6em}{}

\titleformat{\subsubsection}
  {\normalfont\bfseries\itshape}
  {\thesubsubsection}{0.6em}{}

% Thêm dấu chấm trong mục lục cho cấp chapter
\renewcommand{\cftchapdotsep}{\cftdotsep}
\renewcommand{\cftchapfont}{\normalsize\bfseries}

% ============================================================
% 7. THÔNG TIN CÁ NHÂN - THAY PHẦN NÀY
% ============================================================
\newcommand{\tenSV}{Lê Gia Bảo}
\newcommand{\maSV}{<MSSV>}
\newcommand{\tenDoAn}{XÂY DỰNG HỆ THỐNG TRỢ LÝ PHÁP LÝ AI CHO DOANH NGHIỆP NHỎ VÀ VỪA}
\newcommand{\tenNganh}{CÔNG NGHỆ THÔNG TIN}
\newcommand{\namBaoCao}{2026}
\newcommand{\gvhd}{ThS. Dương Hữu Thành}

% Số liệu về corpus và bộ kiểm thử. Khai báo một chỗ để mọi chương dùng chung,
% cập nhật ở đây khi corpus thay đổi.
\newcommand{\soLuat}{38}
\newcommand{\soDieu}{3.844}
\newcommand{\soChuong}{452}
\newcommand{\soRule}{12}
\newcommand{\soTest}{77}

\begin{document}

% Không đánh số trang ở hai trang bìa
\pagenumbering{gobble}

% ============================================================
% TRANG BÌA 1
% ============================================================
\begin{titlepage}
\centering

{\bfseries BỘ GIÁO DỤC VÀ ĐÀO TẠO\par}
\vspace{0.3cm}
{\bfseries TRƯỜNG ĐẠI HỌC MỞ THÀNH PHỐ HỒ CHÍ MINH\par}

\vspace{3cm}

{\large\bfseries \tenSV\par}

\vspace{1cm}

{\Large\bfseries \tenDoAn\par}

\vspace{1.5cm}

{\huge\bfseries ĐỒ ÁN NGÀNH\par}

\vspace{0.8cm}

{\large\bfseries NGÀNH \tenNganh\par}

\vfill

{\bfseries TP. HỒ CHÍ MINH, \namBaoCao\par}

\end{titlepage}

\clearpage

% ============================================================
% TRANG BÌA 2
% ============================================================
\begin{titlepage}
\centering

{\bfseries BỘ GIÁO DỤC VÀ ĐÀO TẠO\par}
\vspace{0.3cm}
{\bfseries TRƯỜNG ĐẠI HỌC MỞ THÀNH PHỐ HỒ CHÍ MINH\par}

\vspace{2.2cm}

{\large\bfseries \tenSV\par}

\vspace{0.5cm}

{\Large\bfseries \tenDoAn\par}

\vspace{0.8cm}

{\large Mã số sinh viên: \maSV\par}

\vspace{1.2cm}

{\huge\bfseries ĐỒ ÁN NGÀNH\par}

\vspace{0.8cm}

{\large\bfseries NGÀNH \tenNganh\par}

\vspace{1.2cm}

{\large Giảng viên hướng dẫn: \gvhd\par}

\vfill

{\bfseries TP. HỒ CHÍ MINH, \namBaoCao\par}

\end{titlepage}

\clearpage

% Bắt đầu đánh số La Mã cho phần đầu báo cáo
\pagenumbering{Roman}
\setcounter{page}{1}
\pagestyle{plain}

% ============================================================
% LỜI CẢM ƠN
% ============================================================
\chapter*{LỜI CẢM ƠN}
% Nếu muốn đưa vào mục lục thì bỏ dấu % dòng sau:
% \addcontentsline{toc}{chapter}{LỜI CẢM ƠN}

Em xin gửi lời cảm ơn chân thành tới ThS.~Dương Hữu Thành, giảng viên hướng dẫn,
đã định hướng đề tài, góp ý về kiến trúc hệ thống và theo sát quá trình hiện thực.
Em cũng cảm ơn các thầy cô Khoa Công nghệ Thông tin, Trường Đại học Mở Thành phố
Hồ Chí Minh, đã trang bị nền tảng để em hoàn thành đồ án này.

Phần lõi truy hồi của hệ thống kế thừa pipeline đã xây dựng trong cuộc thi
Vietnamese Legal AI; đồ án mở rộng phần đó thành sản phẩm ứng dụng hoàn chỉnh cho
doanh nghiệp nhỏ và vừa. Mọi sai sót còn lại thuộc về trách nhiệm của em.

\clearpage

% ============================================================
% NHẬN XÉT CỦA GIÁO VIÊN HƯỚNG DẪN
% ============================================================
\chapter*{NHẬN XÉT CỦA GIÁO VIÊN HƯỚNG DẪN}
% Nếu muốn đưa vào mục lục thì bỏ dấu % dòng sau:
% \addcontentsline{toc}{chapter}{NHẬN XÉT CỦA GIÁO VIÊN HƯỚNG DẪN}

\begin{flushright}
\textit{TP. Hồ Chí Minh, ngày ... tháng ... năm \namBaoCao}\\[2.5cm]
Giảng viên hướng dẫn\\[2.5cm]
\gvhd
\end{flushright}

\clearpage

% ============================================================
% TÓM TẮT ĐỒ ÁN NGÀNH
% ============================================================
\chapter*{TÓM TẮT ĐỒ ÁN NGÀNH}
% Nếu muốn đưa vào mục lục thì bỏ dấu % dòng sau:
% \addcontentsline{toc}{chapter}{TÓM TẮT ĐỒ ÁN NGÀNH}

Doanh nghiệp nhỏ và vừa chiếm phần lớn số lượng doanh nghiệp tại Việt Nam nhưng
hiếm khi có đủ nguồn lực để duy trì bộ phận pháp chế riêng. Trong khi đó, các
nghĩa vụ về thuế, lao động, bảo hiểm xã hội và kế toán được quy định phân tán
trên nhiều văn bản có hiệu lực khác nhau và viện dẫn chéo lẫn nhau, khiến việc
tra cứu thủ công vừa mất thời gian vừa dễ bỏ sót.

Đồ án xây dựng một ứng dụng web trợ lý pháp lý cho nhóm người dùng này, gồm bốn
nghiệp vụ: hỏi đáp pháp luật có trích dẫn tới từng Điều, tra cứu văn bản theo cây
Chương--Điều kèm tìm kiếm toàn văn tiếng Việt, soát xét rủi ro hợp đồng từ tệp
PDF/DOCX, và lịch nhắc nghĩa vụ tuân thủ định kỳ. Lõi hỏi đáp dùng kiến trúc
Retrieval-Augmented Generation với truy hồi kết hợp vector và BM25, trên corpus
gồm \soLuat{} văn bản và \soDieu{} Điều luật liên quan trực tiếp tới doanh nghiệp
nhỏ và vừa. Mỗi câu trả lời đều kèm danh sách Điều luật căn cứ để người dùng tự
kiểm chứng, thay vì buộc phải tin vào kết quả của mô hình.

Hệ thống được hiện thực bằng React và FastAPI trên PostgreSQL, đóng gói bằng
Docker Compose và kiểm thử bằng \soTest{} test tự động.

\clearpage

% ============================================================
% ABSTRACT
% ============================================================
\chapter*{ABSTRACT}
% Nếu muốn đưa vào mục lục thì bỏ dấu % dòng sau:
% \addcontentsline{toc}{chapter}{ABSTRACT}

Small and medium-sized enterprises make up the majority of registered businesses
in Vietnam, yet few can afford an in-house legal department. Their recurring
obligations concerning taxation, labour, social insurance and accounting are
spread across many legal documents of differing authority that cross-reference
one another, which makes manual lookup both slow and error-prone.

This project delivers a web application that assists these businesses across
four workflows: legal question answering with citations down to the individual
article, document browsing over a chapter/article tree with Vietnamese full-text
search, contract risk review from uploaded PDF/DOCX files, and a recurring
compliance calendar. The question answering core follows a
Retrieval-Augmented Generation architecture combining dense vector retrieval with
BM25 over a corpus of \soLuat{} legal documents and \soDieu{} articles selected
for their relevance to small businesses. Every answer carries the list of
articles it relies on, so users can verify the response rather than having to
trust the model.

The system is implemented with React and FastAPI on top of PostgreSQL, packaged
with Docker Compose, and covered by \soTest{} automated tests.

\clearpage

% ============================================================
% MỤC LỤC
% ============================================================
\tableofcontents
\clearpage

% ============================================================
% DANH MỤC TỪ VIẾT TẮT
% ============================================================
\phantomsection
\addcontentsline{toc}{chapter}{DANH MỤC TỪ VIẾT TẮT}
\chapter*{DANH MỤC TỪ VIẾT TẮT}

\begin{longtable}{p{2.8cm}p{5.5cm}p{6cm}}
\toprule
\textbf{Từ viết tắt} & \textbf{Từ đầy đủ} & \textbf{Ý nghĩa} \\
\midrule
\endfirsthead

\toprule
\textbf{Từ viết tắt} & \textbf{Từ đầy đủ} & \textbf{Ý nghĩa} \\
\midrule
\endhead

AI & Artificial Intelligence & Trí tuệ nhân tạo. \\
API & Application Programming Interface & Giao diện lập trình ứng dụng. \\
BHXH & Bảo hiểm xã hội & Nghĩa vụ đóng bảo hiểm của doanh nghiệp. \\
BM25 & Best Matching 25 & Hàm xếp hạng theo từ khóa của truy hồi thưa. \\
CNTT & Công nghệ thông tin & Lĩnh vực công nghệ thông tin. \\
CSDL & Cơ sở dữ liệu & Nơi lưu trữ dữ liệu có cấu trúc. \\
DNNVV & Doanh nghiệp nhỏ và vừa & Nhóm người dùng mục tiêu của hệ thống. \\
FTS & Full-Text Search & Tìm kiếm toàn văn trên văn bản. \\
GTGT & Giá trị gia tăng & Loại thuế khai theo tháng hoặc theo quý. \\
HyDE & Hypothetical Document Embeddings & Sinh câu trả lời giả để mở rộng truy vấn. \\
JWT & JSON Web Token & Chuẩn token dùng cho xác thực không trạng thái. \\
LLM & Large Language Model & Mô hình ngôn ngữ lớn. \\
ORM & Object-Relational Mapping & Ánh xạ đối tượng với bảng quan hệ. \\
RAG & Retrieval-Augmented Generation & Sinh câu trả lời dựa trên ngữ cảnh truy hồi. \\
RRF & Reciprocal Rank Fusion & Kỹ thuật hợp nhất nhiều bảng xếp hạng. \\
SPA & Single Page Application & Ứng dụng web một trang. \\
SSE & Server-Sent Events & Giao thức đẩy dữ liệu một chiều từ server. \\
TNCN & Thu nhập cá nhân & Loại thuế doanh nghiệp khấu trừ cho người lao động. \\
TNDN & Thu nhập doanh nghiệp & Loại thuế tính trên lợi nhuận doanh nghiệp. \\
VBQPPL & Văn bản quy phạm pháp luật & Văn bản chứa quy phạm do cơ quan nhà nước ban hành. \\

\bottomrule
\end{longtable}

\clearpage

% ============================================================
% DANH MỤC HÌNH VẼ
% ============================================================
\phantomsection
\addcontentsline{toc}{chapter}{DANH MỤC HÌNH VẼ}
\listoffigures
\clearpage

% ============================================================
% DANH MỤC BẢNG
% ============================================================
\phantomsection
\addcontentsline{toc}{chapter}{DANH MỤC BẢNG}
\listoftables
\clearpage

% ============================================================
% BẮT ĐẦU NỘI DUNG CHÍNH - ĐÁNH SỐ TRANG Ả RẬP
% ============================================================
\pagenumbering{arabic}
\setcounter{page}{1}

% ============================================================
% CHƯƠNG 1. TỔNG QUAN ĐỀ TÀI
% ============================================================
\chapter{Tổng quan đề tài}

\section{Giới thiệu đề tài}
Đồ án xây dựng một ứng dụng web đóng vai trò trợ lý pháp lý cho doanh nghiệp nhỏ
và vừa (DNNVV) tại Việt Nam. Người dùng mục tiêu là chủ doanh nghiệp, kế toán và
nhân sự --- những người thường xuyên phải tra cứu quy định pháp luật trong công
việc hàng ngày nhưng không được đào tạo về luật.

Bối cảnh của bài toán nằm ở khoảng cách giữa nhu cầu và nguồn lực. Một doanh
nghiệp mười lao động vẫn phải khai thuế giá trị gia tăng, quyết toán thuế thu
nhập doanh nghiệp, đóng bảo hiểm xã hội, báo cáo tình hình sử dụng lao động và
lập báo cáo tài chính năm --- về cơ bản là cùng một tập nghĩa vụ như doanh nghiệp
nghìn lao động, nhưng không có bộ phận pháp chế để theo dõi. Khi cần tra cứu, họ
đứng trước một hệ thống VBQPPL phân tán: quy định gốc nằm trong luật, điều kiện
áp dụng nằm trong nghị định, cách thực hiện nằm trong thông tư, và ba văn bản này
viện dẫn chéo lẫn nhau.

Các công cụ hiện có chưa giải quyết được khoảng cách này. Cổng thông tin VBQPPL
của nhà nước lưu trữ đầy đủ nhưng chỉ tìm theo từ khóa và số hiệu, đòi hỏi người
tra cứu phải biết trước mình cần văn bản nào. Ở đầu ngược lại, các chatbot đa
dụng trả lời được câu hỏi diễn đạt tự nhiên nhưng không nêu căn cứ kiểm chứng
được, và có thể trích dẫn một Điều luật không tồn tại bằng giọng điệu chắc chắn
như khi trích dẫn đúng --- rủi ro không thể chấp nhận trong lĩnh vực mà hệ quả
của thông tin sai là bị xử phạt.

\section{Lý do chọn đề tài}

\subsection{Lý do thực tiễn}
DNNVV chiếm phần lớn số doanh nghiệp đang hoạt động tại Việt Nam nhưng lại là
nhóm ít khả năng chi trả cho tư vấn pháp lý nhất. Hệ quả của việc thiếu thông tin
không nằm ở chỗ tra cứu bất tiện, mà ở chỗ chậm hạn khai thuế hay báo cáo lao
động đều dẫn tới xử phạt hành chính. Một công cụ trả lời được câu hỏi diễn đạt
bằng ngôn ngữ thường ngày, kèm căn cứ pháp lý kiểm chứng được và nhắc trước hạn
của các nghĩa vụ định kỳ, có giá trị sử dụng trực tiếp với nhóm này.

\subsection{Lý do kỹ thuật}
Bài toán này đặt ra ba thách thức kỹ thuật đáng nghiên cứu, cả ba đều có tính đặc
thù chứ không phải áp dụng lại giải pháp có sẵn:

\begin{itemize}
    \item \textbf{Truy hồi trên văn bản pháp luật.} Câu hỏi của người dùng và văn
    bản luật dùng hai lớp từ vựng khác nhau: người dùng hỏi ``sa thải nhân viên
    cần báo trước bao lâu'', còn luật viết ``đơn phương chấm dứt hợp đồng lao
    động''. Tìm theo từ khóa thuần bỏ sót vì không khớp mặt chữ; tìm theo vector
    thuần bỏ sót khi câu hỏi chứa số hiệu văn bản cần khớp chính xác. Đây là lý
    do phải kết hợp cả hai hướng.
    \item \textbf{Xử lý tiếng Việt.} PostgreSQL không có từ điển full-text tiếng
    Việt, và người dùng thường gõ không dấu. Phải tự thiết kế chiến lược chuẩn
    hóa dấu mà không làm mất dấu trong kết quả trả về cho người đọc.
    \item \textbf{Ràng buộc trích dẫn.} Yêu cầu ``mỗi câu trả lời phải kèm Điều
    luật căn cứ và Điều đó phải thật'' biến bài toán sinh văn bản tự do thành bài
    toán sinh có kiểm chứng, đòi hỏi hậu kiểm trích dẫn trên tập văn bản đã truy
    hồi.
\end{itemize}

\section{Phát biểu bài toán}

Bài toán trung tâm được phát biểu như sau. \textbf{Đầu vào} là một câu hỏi pháp
lý diễn đạt bằng tiếng Việt tự nhiên, có thể không dấu, không nêu tên văn bản, do
người không có chuyên môn luật đặt ra. \textbf{Đầu ra} là câu trả lời bằng tiếng
Việt kèm danh sách Điều luật làm căn cứ, trong đó mỗi Điều được dẫn phải tồn tại
trong corpus và liên kết được tới toàn văn để người dùng tự đối chiếu.

Các yêu cầu chính đặt ra cho lời giải:

\begin{itemize}
    \item Trả lời dựa trên văn bản truy hồi được, không dựa vào tri thức ghi
    trong tham số mô hình --- để câu trả lời luôn có căn cứ kiểm chứng và cập
    nhật được khi luật thay đổi.
    \item Không bịa trích dẫn. Khi corpus không chứa quy định liên quan, hệ thống
    phải nói rõ là không tìm thấy thay vì đưa ra câu trả lời chung chung.
    \item Phản hồi tăng dần theo thời gian thực, vì một lượt hỏi đáp có truy hồi
    và sinh văn bản mất hàng chục giây --- nếu chỉ hiện vòng xoay chờ thì người
    dùng sẽ nghĩ hệ thống bị treo.
    \item Cách ly dữ liệu giữa các doanh nghiệp: hợp đồng và lịch tuân thủ là
    thông tin nội bộ.
\end{itemize}

Ngoài bài toán trung tâm, đồ án còn giải ba bài toán phụ trợ: tra cứu văn bản
theo cấu trúc và theo từ khóa, soát xét rủi ro hợp đồng tải lên, và sinh lịch
nghĩa vụ tuân thủ theo đặc điểm doanh nghiệp.

\section{Mục tiêu đề tài}

\textbf{Mục tiêu tổng quát:} xây dựng và triển khai được một ứng dụng web hoàn
chỉnh giúp DNNVV tra cứu, hỏi đáp và theo dõi nghĩa vụ pháp lý, trong đó mọi
thông tin pháp lý đưa ra đều truy vết được về văn bản gốc.

\textbf{Mục tiêu cụ thể:}
\begin{itemize}
    \item Xây dựng corpus pháp luật cho DNNVV: thu thập, chuẩn hóa và chia nhỏ
    theo cấp độ Điều, kèm metadata về loại văn bản, hiệu lực và viện dẫn chéo.
    \item Cài đặt pipeline hỏi đáp RAG với truy hồi kết hợp vector và BM25, có
    bước tinh lọc ngữ cảnh và hậu kiểm trích dẫn.
    \item Xây dựng chức năng tra cứu văn bản theo cây Chương--Điều kèm tìm kiếm
    toàn văn xử lý được truy vấn tiếng Việt không dấu.
    \item Xây dựng chức năng soát xét hợp đồng: trích văn bản từ PDF/DOCX, tách
    theo điều khoản, đối chiếu với quy định pháp luật và chấm mức rủi ro.
    \item Xây dựng lịch tuân thủ sinh tự động theo loại hình và quy mô doanh
    nghiệp, mỗi nghĩa vụ dẫn chiếu về Điều luật quy định nó.
    \item Cài đặt xác thực, phân quyền theo bốn vai trò và cách ly dữ liệu theo
    doanh nghiệp.
    \item Đóng gói hệ thống bằng Docker Compose và kiểm thử tự động cả backend
    lẫn frontend.
\end{itemize}

\section{Phạm vi đề tài}

\subsection{Phạm vi chức năng}
Hệ thống hỗ trợ bốn nghiệp vụ: hỏi đáp pháp luật, tra cứu văn bản, soát xét hợp
đồng và lịch tuân thủ. Hệ thống \textbf{không} thay thế ý kiến của luật sư, không
đại diện doanh nghiệp trong tố tụng và không tự động kê khai với cơ quan nhà
nước. Lịch tuân thủ nhắc hạn và dẫn căn cứ pháp lý nhưng không nộp hồ sơ thay
người dùng.

\subsection{Phạm vi kỹ thuật}
Sản phẩm là ứng dụng web chạy trên trình duyệt máy tính, chưa có ứng dụng di động
riêng. Hệ thống gọi LLM và embedding qua API của nhà cung cấp bên ngoài thay vì
tự triển khai mô hình, do hạn chế về hạ tầng tính toán. Xử lý nền dùng
\texttt{BackgroundTasks} của FastAPI, phù hợp quy mô đồ án; hệ thống chưa có hàng
đợi tác vụ phân tán.

\subsection{Phạm vi dữ liệu}
Corpus gồm \soLuat{} văn bản với \soDieu{} Điều, chọn theo tiêu chí liên quan
trực tiếp tới hoạt động của DNNVV: doanh nghiệp, thuế, lao động, bảo hiểm xã hội,
kế toán, hợp đồng dân sự, sở hữu trí tuệ và hỗ trợ DNNVV. Đây là dữ liệu thật lấy
từ VBQPPL đã công bố, không phải dữ liệu mô phỏng. Corpus không bao gồm luật hình
sự, luật hành chính chuyên ngành hay pháp luật địa phương, và phản ánh hiệu lực
tại thời điểm thu thập --- văn bản sửa đổi ban hành sau đó cần được nhập bổ sung
qua chức năng quản trị.

\section{Phương pháp thực hiện đề tài}

\textbf{Khảo sát và thu thập yêu cầu.} Yêu cầu chức năng được suy ra từ các nghĩa
vụ pháp lý định kỳ có thật của DNNVV, tra cứu trực tiếp trong VBQPPL về thuế, lao
động và kế toán. Cách làm này cho ra một tập yêu cầu kiểm chứng được: mỗi nghĩa vụ
trong lịch tuân thủ đều dẫn được về Điều luật quy định nó.

\textbf{Phân tích và thiết kế.} Hệ thống được mô hình hóa bằng UML: sơ đồ use case
xác định phạm vi chức năng theo từng vai trò, sơ đồ lớp thực thể làm cơ sở cho
lược đồ CSDL, sơ đồ tuần tự đặc tả các luồng cốt lõi, sơ đồ trạng thái mô tả vòng
đời tác vụ soát xét hợp đồng.

\textbf{Phát triển.} Quá trình phát triển đi theo hướng tăng trưởng dần: dựng nền
tảng CSDL và xác thực trước, sau đó thêm từng nghiệp vụ, mỗi nghiệp vụ hoàn thiện
cả API lẫn giao diện trước khi chuyển sang nghiệp vụ tiếp theo. Cách này giúp
luôn có một hệ thống chạy được ở mọi thời điểm.

\textbf{Cài đặt và triển khai.} Toàn bộ hệ thống được đóng gói bằng Docker
Compose gồm ba container, khởi động bằng một lệnh duy nhất, migration CSDL chạy
tự động lúc container backend khởi động.

\textbf{Kiểm thử và đánh giá.} Backend kiểm thử bằng \texttt{pytest} trên một
database riêng tạo tự động rồi hủy sau khi chạy; frontend kiểm thử bằng
\texttt{Vitest}. Chất lượng truy hồi được đo bằng độ chính xác, độ phủ và điểm
F2 trên bộ câu hỏi có đáp án Điều luật.

\section{Khảo sát liên quan hệ thống}

Bảng~\ref{tab:khao-sat} so sánh ba nhóm giải pháp mà người dùng DNNVV hiện có thể
tiếp cận.

\begin{table}[htp]
    \centering
    \caption{So sánh các giải pháp tra cứu pháp luật hiện có.}
    \label{tab:khao-sat}
    \begin{tabular}{p{3.1cm}p{3.5cm}p{3.5cm}p{3.5cm}}
        \toprule
        \textbf{Tiêu chí} & \textbf{Cổng VBQPPL} & \textbf{Chatbot đa dụng} &
        \textbf{Hệ thống của đồ án} \\
        \midrule
        Cách tra cứu & Từ khóa, số hiệu & Câu hỏi tự nhiên & Câu hỏi tự nhiên
        và tra cứu theo cấu trúc \\
        Căn cứ pháp lý & Chính là văn bản gốc & Không kiểm chứng được & Trích dẫn
        Điều, liên kết tới toàn văn \\
        Nguy cơ bịa & Không có & Cao & Hạn chế bằng hậu kiểm trên corpus \\
        Nghiệp vụ SME & Không & Không & Soát xét hợp đồng, lịch tuân thủ \\
        Phạm vi & Toàn bộ VBQPPL & Không xác định & \soLuat{} văn bản liên quan
        DNNVV \\
        \bottomrule
    \end{tabular}
\end{table}

Khoảng trống mà đồ án hướng tới nằm ở chỗ: cổng VBQPPL đáng tin nhưng khó dùng
với người không biết mình cần văn bản nào, còn chatbot đa dụng dễ dùng nhưng
không đáng tin. Hệ thống của đồ án cố gắng giữ tính dễ dùng của hướng thứ hai
trong khi bù lại độ tin cậy bằng cách buộc mọi câu trả lời phải neo vào một tập
văn bản xác định và kiểm chứng được.

Về mặt kỹ thuật, đồ án kế thừa hai kết quả đã được kiểm chứng trong tài liệu:
kiến trúc RAG với truy hồi kết hợp dense và sparse cho bài toán hỏi đáp có căn cứ
tài liệu~\cite{lewis2020rag, karpukhin2020dpr}, và kỹ thuật hợp nhất xếp
hạng~\cite{cormack2009rrf}.

\section{Bố cục báo cáo}

\begin{itemize}
    \item \textbf{Chương 1. Tổng quan đề tài:} bối cảnh pháp lý của DNNVV, lý do
    chọn đề tài, phát biểu bài toán hỏi đáp pháp luật có trích dẫn, mục tiêu,
    phạm vi và phương pháp thực hiện.
    \item \textbf{Chương 2. Cơ sở lý thuyết:} cấu trúc hệ thống VBQPPL Việt Nam,
    các kỹ thuật truy hồi thông tin (BM25, dense retrieval, RRF), kiến trúc RAG
    và các công nghệ được dùng trong đồ án.
    \item \textbf{Chương 3. Hệ thống:} kiến trúc ba lớp, đặc tả use case, mô hình
    hóa UML, thiết kế CSDL, thiết kế giao diện và kết quả triển khai kèm đánh giá.
    \item \textbf{Chương 4. Kết luận và hướng phát triển:} tổng kết kết quả, hạn
    chế còn tồn tại và các hướng mở rộng.
\end{itemize}


% ============================================================
% CHƯƠNG 2. CƠ SỞ LÝ THUYẾT
% ============================================================
\chapter{Cơ sở lý thuyết}

\section{Cấu trúc hệ thống văn bản quy phạm pháp luật Việt Nam}

\subsection{Phân cấp hiệu lực pháp lý}
VBQPPL Việt Nam được tổ chức theo thứ bậc hiệu lực. Trong phạm vi đề tài, ba cấp
quan trọng nhất là:

\begin{itemize}
    \item \textbf{Luật} và \textbf{Bộ luật} do Quốc hội ban hành, quy định các
    nguyên tắc và quyền, nghĩa vụ cơ bản.
    \item \textbf{Nghị định} do Chính phủ ban hành, quy định chi tiết các điều
    khoản mà luật giao, thường là nơi chứa điều kiện áp dụng và mức xử phạt.
    \item \textbf{Thông tư} do Bộ trưởng ban hành, hướng dẫn cách thực hiện: biểu
    mẫu, trình tự, thời hạn.
\end{itemize}

Thứ bậc này có hệ quả trực tiếp lên thiết kế hệ thống. Câu trả lời đầy đủ cho một
câu hỏi thực tế thường phải tổng hợp cả ba cấp: luật cho biết nghĩa vụ tồn tại,
nghị định cho biết khi nào áp dụng, thông tư cho biết làm thế nào. Vì vậy hệ thống
không thể chỉ truy hồi một Điều rồi trả lời, mà phải lấy được một tập Điều thuộc
nhiều văn bản, và corpus buộc phải chứa đủ cả ba cấp --- đây là lý do corpus của
đồ án gồm cả Luật, Nghị định và Thông tư như trình bày ở mục~\ref{sec:corpus}.

\subsection{Cấu trúc nội tại của một văn bản}
Một VBQPPL được chia theo cấp: Phần $\rightarrow$ Chương $\rightarrow$ Mục
$\rightarrow$ Điều $\rightarrow$ Khoản $\rightarrow$ Điểm. Trong đó \textbf{Điều}
là đơn vị đặc biệt quan trọng: nó là đơn vị nhỏ nhất có tiêu đề riêng, diễn đạt
một quy phạm tương đối trọn nghĩa, và là đơn vị mà giới hành nghề dùng để dẫn
chiếu. Nhận xét này là cơ sở cho chiến lược chia nhỏ dữ liệu ở
mục~\ref{sec:chunking}.

\subsection{Hiện tượng viện dẫn chéo}
Các Điều luật thường xuyên dẫn chiếu lẫn nhau, cả trong cùng văn bản (``theo quy
định tại Điều 35 của Luật này'') lẫn sang văn bản khác. Một Điều bị tách khỏi các
Điều mà nó dẫn chiếu có thể mất nghĩa hoặc bị hiểu sai. Hệ thống vì vậy cần lưu
lại quan hệ viện dẫn dưới dạng dữ liệu có cấu trúc, để hiển thị cho người đọc
những Điều liên quan thay vì để họ tự đi tìm.

\section{Kỹ thuật truy hồi thông tin}

\subsection{Truy hồi thưa và thuật toán BM25}
Truy hồi thưa (\textit{sparse retrieval}) biểu diễn văn bản bằng vector trên không
gian từ vựng, mỗi chiều ứng với một từ. BM25~\cite{robertson2009bm25} là hàm xếp
hạng tiêu biểu, tính điểm liên quan giữa truy vấn $Q$ và tài liệu $D$ theo công
thức:

\begin{equation}
\label{eq:bm25}
\mathrm{BM25}(Q, D) = \sum_{t \in Q} \mathrm{IDF}(t) \cdot
\frac{f(t, D) \cdot (k_1 + 1)}
     {f(t, D) + k_1 \cdot \left(1 - b + b \cdot \dfrac{|D|}{\mathrm{avgdl}}\right)}
\end{equation}

trong đó $f(t, D)$ là số lần từ $t$ xuất hiện trong $D$, $|D|$ là độ dài tài
liệu, $\mathrm{avgdl}$ là độ dài trung bình của tài liệu trong tập, còn $k_1$ và
$b$ là siêu tham số điều khiển mức bão hòa tần số từ và mức chuẩn hóa theo độ
dài. Thành phần $\mathrm{IDF}(t)$ hạ trọng số của những từ xuất hiện ở khắp nơi:

\begin{equation}
\label{eq:idf}
\mathrm{IDF}(t) = \ln\left(\frac{N - n(t) + 0.5}{n(t) + 0.5} + 1\right)
\end{equation}

với $N$ là tổng số tài liệu và $n(t)$ là số tài liệu chứa $t$.

Điểm mạnh của BM25 trong bài toán pháp lý là khớp chính xác: khi người dùng gõ
``Điều 35 Bộ luật Lao động'' hay một số hiệu văn bản, BM25 tìm đúng, còn truy hồi
theo vector có thể trả về một Điều gần nghĩa nhưng sai số. Điểm yếu là nó hoàn
toàn không xử lý được khác biệt từ vựng: ``sa thải'' và ``đơn phương chấm dứt hợp
đồng lao động'' không có từ nào chung nên điểm BM25 bằng không.

Áp dụng BM25 cho tiếng Việt cần một bước tiền xử lý riêng. Tiếng Việt không phân
tách từ bằng khoảng trắng --- ``hợp đồng lao động'' là ba âm tiết tạo thành một
thuật ngữ --- nên cần tách từ bằng công cụ chuyên biệt cho tiếng Việt trước khi
tính tần số.

\subsection{Truy hồi dày đặc}
Truy hồi dày đặc (\textit{dense retrieval})~\cite{karpukhin2020dpr} biểu diễn cả
truy vấn và tài liệu bằng vector số thực trong không gian ngữ nghĩa nhiều chiều,
sinh ra bởi một mô hình embedding. Độ liên quan được đo bằng độ tương đồng
cosine:

\begin{equation}
\label{eq:cosine}
\mathrm{sim}(q, d) = \frac{\mathbf{v}_q \cdot \mathbf{v}_d}
                          {\lVert \mathbf{v}_q \rVert \, \lVert \mathbf{v}_d \rVert}
\end{equation}

Vì vector được huấn luyện để hai đoạn văn cùng nghĩa nằm gần nhau bất kể dùng từ
nào, hướng này giải quyết được đúng điểm yếu của BM25. Đánh đổi là nó mờ với
những chi tiết cần khớp chính xác: mô hình embedding coi ``Điều 35'' và ``Điều
36'' là gần như đồng nhất.

Hai điểm mạnh và điểm yếu này gần như bù trừ nhau, và đó là lập luận để dùng cả
hai đồng thời.

\subsection{Truy hồi kết hợp và Reciprocal Rank Fusion}
Truy hồi kết hợp (\textit{hybrid retrieval}) chạy song song hai hướng trên rồi
hợp nhất kết quả. Trở ngại là điểm số của hai hướng không cùng thang: BM25 cho
điểm không chặn trên, cosine cho điểm trong $[-1, 1]$, nên không thể cộng trực
tiếp.

Reciprocal Rank Fusion~\cite{cormack2009rrf} giải quyết bằng cách bỏ hẳn điểm số
và chỉ dùng thứ hạng:

\begin{equation}
\label{eq:rrf}
\mathrm{RRF}(d) = \sum_{i} \frac{w_i}{k + \mathrm{rank}_i(d)}
\end{equation}

trong đó $\mathrm{rank}_i(d)$ là thứ hạng của tài liệu $d$ trong danh sách thứ
$i$, $w_i$ là trọng số của danh sách đó, và $k$ là hằng số làm phẳng (thường
chọn $k = 60$) nhằm giảm ảnh hưởng quá lớn của vài vị trí đầu. Cách này bền vững
vì chỉ phụ thuộc thứ tự tương đối, không phụ thuộc thang điểm; trọng số $w_i$ cho
phép ưu tiên hướng truy hồi tỏ ra hiệu quả hơn trên miền dữ liệu cụ thể.

\subsection{Chỉ số đánh giá truy hồi}
Với một truy vấn có tập Điều đúng $R$ và tập Điều hệ thống trả về $A$, độ chính
xác và độ phủ được định nghĩa:

\begin{equation}
\label{eq:precision-recall}
P = \frac{|R \cap A|}{|A|}, \qquad
R_{\mathrm{ecall}} = \frac{|R \cap A|}{|R|}
\end{equation}

Hai chỉ số này được tổng hợp bằng $F_\beta$:

\begin{equation}
\label{eq:fbeta}
F_\beta = (1 + \beta^2) \cdot
\frac{P \cdot R_{\mathrm{ecall}}}{\beta^2 \cdot P + R_{\mathrm{ecall}}}
\end{equation}

Đồ án chọn $\beta = 2$, tức $F_2$ --- coi độ phủ quan trọng gấp bốn lần độ chính
xác. Lựa chọn này xuất phát từ đặc thù bài toán: nếu truy hồi thừa một Điều không
liên quan, bước tinh lọc phía sau còn cơ hội loại nó ra; nhưng nếu bỏ sót Điều
quyết định, không bước nào sau đó cứu được và câu trả lời sẽ sai.

\section{Kiến trúc Retrieval-Augmented Generation}

\subsection{Nguyên lý}
RAG~\cite{lewis2020rag} kết hợp một thành phần truy hồi với một LLM: thay vì để
mô hình trả lời bằng tri thức ghi trong tham số, hệ thống truy hồi các đoạn văn
bản liên quan rồi đưa vào ngữ cảnh của prompt và yêu cầu mô hình trả lời dựa trên
đó.

Với bài toán pháp lý, RAG mang lại ba lợi ích thiết yếu. Thứ nhất, câu trả lời
neo vào văn bản xác định nên nêu được căn cứ kiểm chứng. Thứ hai, cập nhật luật
chỉ cần nhập văn bản mới vào corpus, không phải huấn luyện lại mô hình. Thứ ba,
có thể hậu kiểm: đối chiếu trích dẫn trong câu trả lời với tập văn bản đã truy
hồi để phát hiện dẫn chiếu bịa.

\subsection{Các rủi ro đặc thù}
RAG không tự động loại bỏ mọi vấn đề. Ba rủi ro cần xử lý bằng thiết kế:

\begin{itemize}
    \item \textbf{Hallucination.} Mô hình vẫn có thể bổ sung thông tin không có
    trong ngữ cảnh. Cách hạn chế là ràng buộc bằng prompt và hậu kiểm trích dẫn.
    \item \textbf{Citation misalignment.} Câu trả lời đúng nhưng dẫn sai số Điều.
    Với người dùng, dạng lỗi này còn nguy hiểm hơn trả lời sai hẳn, vì nó tạo cảm
    giác đã được kiểm chứng.
    \item \textbf{Nhiễu ngữ cảnh.} Đưa quá nhiều đoạn văn bản vào prompt làm loãng
    thông tin quan trọng và hạ chất lượng câu trả lời, nên cần một bước tinh lọc
    giữa truy hồi và sinh.
\end{itemize}

\subsection{Viết lại truy vấn và HyDE}
Câu hỏi của người dùng thường ngắn và dùng từ vựng thông thường, khác xa văn
phong của văn bản luật. Hai kỹ thuật thu hẹp khoảng cách này:

\begin{itemize}
    \item \textbf{Viết lại truy vấn:} dùng LLM diễn đạt lại câu hỏi bằng thuật
    ngữ pháp lý trước khi truy hồi.
    \item \textbf{HyDE}~\cite{gao2023hyde} (\textit{Hypothetical Document
    Embeddings}): yêu cầu LLM sinh một câu trả lời giả định, rồi dùng chính văn
    bản giả định đó làm truy vấn cho truy hồi dày đặc. Ý tưởng là câu trả lời giả
    định --- dù có thể sai về nội dung --- vẫn nằm gần văn bản luật thật trong
    không gian ngữ nghĩa hơn so với câu hỏi gốc.
\end{itemize}

\section{Công nghệ sử dụng}

\subsection{Tổng quan các công nghệ}
Bảng~\ref{tab:cong-nghe} liệt kê các công nghệ chính và vai trò của từng công
nghệ trong hệ thống.

\begin{table}[htp]
    \centering
    \caption{Các công nghệ sử dụng và vai trò trong hệ thống.}
    \label{tab:cong-nghe}
    \begin{tabular}{p{3.4cm}p{4.2cm}p{6.2cm}}
        \toprule
        \textbf{Công nghệ} & \textbf{Loại} & \textbf{Vai trò trong đồ án} \\
        \midrule
        React 19, TypeScript & Thư viện giao diện & Xây dựng SPA cho toàn bộ
        giao diện người dùng \\
        Vite & Công cụ build & Máy chủ phát triển và đóng gói bản phát hành \\
        Tailwind CSS & Framework CSS & Định dạng giao diện theo tiện ích \\
        TanStack Query & Quản lý trạng thái server & Cache, làm mới và đồng bộ
        dữ liệu từ API \\
        FastAPI & Web framework & Toàn bộ REST API và luồng SSE \\
        Pydantic & Kiểm định dữ liệu & Kiểm tra dữ liệu vào/ra và đọc cấu hình \\
        SQLAlchemy 2.0 & ORM & Truy cập CSDL bất đồng bộ \\
        Alembic & Công cụ migration & Quản lý phiên bản lược đồ CSDL \\
        PostgreSQL 17 & Hệ quản trị CSDL & Dữ liệu nghiệp vụ, corpus và tìm kiếm
        toàn văn \\
        ChromaDB & CSDL vector & Lưu và tìm kiếm vector embedding \\
        LangGraph & Framework điều phối & Mô hình hóa pipeline RAG thành đồ thị \\
        Docker Compose & Đóng gói & Triển khai ba container bằng một lệnh \\
        \bottomrule
    \end{tabular}
\end{table}

\subsection{Lý do lựa chọn công nghệ}
Các lựa chọn quan trọng và lập luận kèm theo:

\textbf{FastAPI cho backend.} Yếu tố quyết định là hỗ trợ bất đồng bộ nguyên bản.
Nghiệp vụ của hệ thống chủ yếu là chờ vào/ra --- chờ LLM, chờ embedding, chờ CSDL
--- nên mô hình bất đồng bộ cho phép một tiến trình phục vụ nhiều yêu cầu đồng
thời mà không cần nhiều luồng. FastAPI cũng hỗ trợ trả về luồng dữ liệu, điều
kiện cần cho SSE, và sinh tài liệu OpenAPI tự động từ kiểu dữ liệu Pydantic.

\textbf{React với Vite thay vì framework render phía server.} Ứng dụng nằm sau
đăng nhập nên không cần tối ưu SEO hay render phía server --- những lợi ích chính
mà framework loại đó mang lại. Ngược lại, SPA thuần đơn giản hơn khi phải đọc
luồng SSE và giữ trạng thái phiên chat ở phía client. Chọn SPA vì vậy là loại bỏ
độ phức tạp không cần thiết.

\textbf{PostgreSQL kèm ChromaDB thay vì một CSDL duy nhất.} Hệ thống có hai nhu
cầu truy vấn khác bản chất: dữ liệu nghiệp vụ cần quan hệ, giao dịch và ràng buộc
toàn vẹn --- đúng thế mạnh của PostgreSQL; còn tìm kiếm vector cần chỉ mục xấp xỉ
lân cận gần nhất, là thế mạnh của CSDL vector. Dùng hai hệ chuyên biệt cho hai
nhu cầu tách biệt, thay vì buộc một hệ làm cả hai việc.

\textbf{LangGraph để điều phối pipeline.} Pipeline hỏi đáp gồm nhiều bước có
nhánh điều kiện và cần báo tiến độ ra ngoài sau mỗi bước. Mô hình hóa thành đồ thị
với trạng thái tường minh giúp việc phát tín hiệu tiến độ cho SSE trở thành hệ quả
tự nhiên của cấu trúc, thay vì phải cài thêm cơ chế riêng.

\section{Các nghiên cứu, hệ thống liên quan}

Các công trình nền tảng mà đồ án dựa trên đã trình bày ở các mục trên:
RAG~\cite{lewis2020rag} cho kiến trúc tổng thể, DPR~\cite{karpukhin2020dpr} cho
truy hồi dày đặc, BM25~\cite{robertson2009bm25} cho truy hồi thưa,
RRF~\cite{cormack2009rrf} cho hợp nhất xếp hạng và HyDE~\cite{gao2023hyde} cho mở
rộng truy vấn. Nghiên cứu về ứng dụng RAG cho chatbot hỏi đáp trong miền hẹp tiếng
Việt~\cite{duong2026} cho thấy hướng tiếp cận này khả thi với ngôn ngữ và quy mô
dữ liệu tương đương đề tài.

Bài học rút ra và áp dụng vào đồ án: các hệ thống hỏi đáp trong miền chuyên môn
đạt chất lượng cao không nhờ mô hình sinh mạnh hơn, mà nhờ chất lượng truy hồi và
việc chia nhỏ dữ liệu đúng cấu trúc của miền. Đây là lý do đồ án đầu tư vào chiến
lược chia theo Điều và truy hồi kết hợp, thay vì tập trung tinh chỉnh mô hình
sinh.


% ============================================================
% CHƯƠNG 3. HỆ THỐNG
% ============================================================
\chapter{Hệ thống}

\section{Giới thiệu hệ thống}

Hệ thống là một ứng dụng web đa người dùng, phục vụ nhiều doanh nghiệp trên cùng
một cài đặt với dữ liệu cách ly theo doanh nghiệp. Người dùng đăng ký tài khoản
kèm thông tin doanh nghiệp, sau đó truy cập bốn nghiệp vụ chính.

\textbf{Hỏi đáp pháp luật} là nghiệp vụ trung tâm. Người dùng đặt câu hỏi bằng
tiếng Việt tự nhiên và nhận câu trả lời hiện dần theo thời gian thực, kèm danh
sách Điều luật căn cứ ở dạng liên kết bấm được. Trong lúc chờ, giao diện hiển thị
tiến độ từng bước xử lý để người dùng biết hệ thống đang làm gì.

\textbf{Tra cứu văn bản} cho phép duyệt corpus theo cây Chương--Điều, tìm kiếm
toàn văn có xử lý truy vấn không dấu, xem Điều liên quan qua quan hệ viện dẫn, và
chuyển thẳng một Điều đang đọc sang khung chat để hỏi thêm.

\textbf{Soát xét hợp đồng} nhận tệp PDF hoặc DOCX, tách hợp đồng theo điều khoản,
đối chiếu từng điều khoản với quy định pháp luật và trả về danh sách phát hiện có
xếp mức rủi ro kèm đề xuất chỉnh sửa, mỗi phát hiện dẫn chiếu về Điều luật căn cứ.

\textbf{Lịch tuân thủ} sinh danh sách nghĩa vụ định kỳ theo loại hình và quy mô
doanh nghiệp, phân nhóm theo mức độ cấp thiết của hạn chót và cho phép đánh dấu
hoàn thành.

Ngoài ra hệ thống có khu vực quản trị để nhập bổ sung văn bản, dựng lại chỉ mục
vector và quản lý người dùng.

\section{Kiến trúc hệ thống}

\subsection{Kiến trúc tổng thể hệ thống}

Hệ thống được tổ chức theo kiến trúc ba lớp, mỗi lớp chạy trong một container
riêng như minh họa ở Hình~\ref{fig:kien-truc-he-thong}.

\begin{figure}[htp]
    \centering
    \includegraphics[width=0.75\linewidth]{kien-truc-he-thong.png}
    \caption{Kiến trúc tổng thể hệ thống.}
    \label{fig:kien-truc-he-thong}
\end{figure}

Lớp trình bày là một SPA React được nginx phục vụ dưới dạng tệp tĩnh. Điểm đáng
lưu ý trong thiết kế: nginx đồng thời chuyển tiếp mọi yêu cầu tới đường dẫn
\texttt{/api} sang container backend, nhờ đó trình duyệt thấy frontend và backend
cùng một origin. Cách này loại bỏ hoàn toàn nhu cầu cấu hình CORS ở môi trường
triển khai và tránh được lớp phức tạp đi kèm.

Lớp ứng dụng là backend FastAPI, chia thành các router theo nghiệp vụ. Toàn bộ
logic RAG nằm trong lớp dịch vụ, tách khỏi lớp HTTP --- lựa chọn này cho phép
nghiệp vụ soát xét hợp đồng gọi lại đúng pipeline truy hồi của nghiệp vụ hỏi đáp
mà không phải viết lại.

Lớp dữ liệu gồm PostgreSQL cho dữ liệu nghiệp vụ và corpus, và ChromaDB cho
vector embedding. Chỉ mục BM25 không được lưu ở CSDL mà dựng trong bộ nhớ tiến
trình lúc khởi động, kèm cache trên đĩa gắn với dấu vân tay của corpus: nếu corpus
không đổi thì lần khởi động sau nạp lại cache thay vì tách từ lại toàn bộ
\soDieu{} Điều.

\subsection{Các thành phần của kiến trúc}

Bảng~\ref{tab:thanh-phan} liệt kê các nhóm thành phần phía backend và trách nhiệm
của từng nhóm.

\begin{table}[htp]
    \centering
    \caption{Các thành phần chính của backend.}
    \label{tab:thanh-phan}
    \begin{tabular}{p{3.6cm}p{4.0cm}p{6.2cm}}
        \toprule
        \textbf{Nhóm} & \textbf{Thành phần} & \textbf{Trách nhiệm} \\
        \midrule
        \multirow{5}{*}{Router (HTTP)}
          & \texttt{auth} & Đăng ký, đăng nhập, làm mới token, đổi mật khẩu \\
          & \texttt{legal} & Hỏi đáp: một lượt và luồng SSE \\
          & \texttt{laws} & Tra cứu, cây Chương--Điều, tìm kiếm toàn văn \\
          & \texttt{documents} & Tải lên và soát xét hợp đồng \\
          & \texttt{compliance} & Nghĩa vụ tuân thủ \\
        \midrule
        \multirow{3}{*}{Dịch vụ lõi}
          & \texttt{agents} & Pipeline RAG dạng đồ thị \\
          & \texttt{vector\_store} & Truy hồi kết hợp và dựng chỉ mục \\
          & \texttt{legal} & Tra cứu và tìm kiếm toàn văn \\
        \midrule
        \multirow{2}{*}{Hạ tầng}
          & \texttt{core.security} & Băm mật khẩu, phát và kiểm JWT \\
          & \texttt{core.database} & Engine và session bất đồng bộ \\
        \bottomrule
    \end{tabular}
\end{table}

\subsection{Luồng xử lý tổng quát của hệ thống}
\label{sec:pipeline}

Luồng hỏi đáp là luồng phức tạp nhất, được mô hình hóa thành đồ thị bảy bước:

\begin{enumerate}
    \item \textbf{Phân tích ý định.} Xác định câu hỏi có thuộc miền pháp luật hay
    không. Câu hỏi ngoài miền được trả lời ngay, không tốn chi phí truy hồi.
    \item \textbf{Chuẩn bị truy vấn.} Viết lại câu hỏi bằng thuật ngữ pháp lý và
    áp dụng HyDE để thu hẹp khoảng cách từ vựng giữa câu hỏi và văn bản luật.
    \item \textbf{Truy hồi.} Chạy song song tìm kiếm vector và BM25, hợp nhất
    bằng RRF có trọng số theo công thức~(\ref{eq:rrf}).
    \item \textbf{Xếp hạng lại.} Bước tùy chọn, dùng mô hình cross-encoder khi có
    endpoint tương ứng.
    \item \textbf{Tinh lọc bằng LLM.} Loại các Điều được truy hồi nhưng không
    thực sự liên quan, nhằm giảm nhiễu ngữ cảnh.
    \item \textbf{Sinh câu trả lời.} Sinh câu trả lời theo từng token dựa trên
    tập Điều đã tinh lọc, phát ra ngoài ngay khi có.
    \item \textbf{Chuẩn hóa kết quả.} Trích danh sách Điều được dẫn và đối chiếu
    với tập đã truy hồi để loại dẫn chiếu không có căn cứ.
\end{enumerate}

Sau mỗi bước, đồ thị phát một sự kiện tiến độ ra luồng SSE. Nhờ trạng thái của đồ
thị là tường minh, việc báo tiến độ không cần cơ chế riêng mà là hệ quả trực tiếp
của cấu trúc pipeline.

Bước thứ bảy chính là cơ chế chống bịa trích dẫn: nếu mô hình sinh ra một dẫn
chiếu tới Điều không nằm trong tập đã truy hồi, dẫn chiếu đó bị loại khỏi danh
sách căn cứ trả về cho người dùng.

\section{Xây dựng corpus pháp luật}
\label{sec:corpus}

\subsection{Thu thập và chuẩn hóa}
Corpus được xây dựng từ một tệp kê khai liệt kê các văn bản cần thu thập kèm phân
loại lĩnh vực. Tiêu chí chọn văn bản là mức liên quan trực tiếp tới hoạt động
thường ngày của DNNVV. Kết quả gồm \soLuat{} văn bản, \soDieu{} Điều thuộc
\soChuong{} Chương, phân bố như ở Bảng~\ref{tab:corpus}.

\begin{table}[htp]
    \centering
    \caption{Phân bố corpus theo loại văn bản và lĩnh vực.}
    \label{tab:corpus}
    \begin{tabular}{p{3.6cm}p{1.8cm}p{4.6cm}p{1.8cm}}
        \toprule
        \textbf{Loại văn bản} & \textbf{Số VB} & \textbf{Lĩnh vực} &
        \textbf{Số VB} \\
        \midrule
        Luật & 19 & Thuế & 12 \\
        Nghị định & 14 & Lao động & 4 \\
        Thông tư & 3 & Sở hữu trí tuệ & 4 \\
        Bộ luật & 2 & Doanh nghiệp & 3 \\
         & & Dân sự -- Hợp đồng & 3 \\
         & & Kế toán -- Tài chính & 3 \\
         & & Các lĩnh vực khác & 9 \\
        \midrule
        \textbf{Tổng} & \textbf{38} & \textbf{Tổng} & \textbf{38} \\
        \bottomrule
    \end{tabular}
\end{table}

Phân bố này phản ánh trọng tâm nghiệp vụ: thuế chiếm tỷ trọng lớn nhất vì đây là
lĩnh vực có nhiều nghĩa vụ định kỳ nhất và cũng là nơi doanh nghiệp dễ vi phạm
nhất do sai sót thủ tục.

\subsection{Chiến lược chia nhỏ dữ liệu}
\label{sec:chunking}

Hệ thống chia corpus theo cấp độ \textbf{Điều}: mỗi Điều luật là một đơn vị lưu
trữ và truy hồi độc lập, độ dài trung bình khoảng 1.600 ký tự.

Quyết định này được lập luận từ ba phương án khả dĩ. Chia theo \textbf{cả văn bản}
tạo ra đơn vị quá lớn --- một bộ luật hàng trăm Điều không thể biểu diễn bằng một
vector duy nhất mà giữ được thông tin, và cũng vượt giới hạn ngữ cảnh khi đưa vào
prompt. Chia theo \textbf{độ dài cố định} là cách thông dụng nhưng ở đây gây hại
trực tiếp: cắt theo số ký tự sẽ chia đôi một Điều ở giữa câu, làm đoạn văn bản mất
nghĩa pháp lý và khiến hệ thống không dẫn chiếu chính xác được. Chia theo
\textbf{Điều} vừa khớp với đơn vị mà giới hành nghề dùng để dẫn chiếu, vừa đủ nhỏ
để nhiều Điều cùng nằm trong ngữ cảnh, và bảo toàn tính trọn nghĩa của quy phạm.

Mỗi Điều được lưu kèm metadata: số hiệu văn bản, tên văn bản, loại văn bản, Chương
chứa nó, tiêu đề Điều, và danh sách viện dẫn chéo. Metadata này phục vụ cả việc
hiển thị nguồn cho người dùng lẫn việc lọc khi truy hồi.

\section{Phân tích sơ đồ use case hệ thống}

\subsection{Sơ đồ use case}

Hệ thống có bốn tác nhân, ứng với bốn vai trò người dùng, phân quyền như ở
Bảng~\ref{tab:vai-tro}.

\begin{table}[htp]
    \centering
    \caption{Các tác nhân và quyền tương ứng.}
    \label{tab:vai-tro}
    \begin{tabular}{p{2.6cm}p{5.2cm}p{6.0cm}}
        \toprule
        \textbf{Tác nhân} & \textbf{Mô tả} & \textbf{Quyền chính} \\
        \midrule
        Chủ doanh nghiệp & Người quản lý doanh nghiệp & Toàn bộ nghiệp vụ trong
        phạm vi doanh nghiệp mình \\
        Kế toán & Phụ trách thuế, kế toán & Hỏi đáp, tra cứu, soát xét hợp đồng,
        lịch tuân thủ \\
        Nhân sự & Phụ trách lao động, BHXH & Hỏi đáp, tra cứu, soát xét hợp đồng,
        lịch tuân thủ \\
        Quản trị viên & Vận hành hệ thống & Thêm quyền quản lý corpus, người dùng
        và chạy đánh giá \\
        \bottomrule
    \end{tabular}
\end{table}

\begin{figure}[htp]
    \centering
    \includegraphics[width=0.85\linewidth]{so-do-use-case.png}
    \caption{Sơ đồ use case của hệ thống.}
    \label{fig:so-do-use-case}
\end{figure}

Danh sách use case chính được tổng hợp ở Bảng~\ref{tab:danh-sach-uc}.

\begin{table}[htp]
    \centering
    \caption{Danh sách use case chính.}
    \label{tab:danh-sach-uc}
    \begin{tabular}{p{1.6cm}p{5.4cm}p{6.6cm}}
        \toprule
        \textbf{Mã} & \textbf{Tên use case} & \textbf{Tác nhân} \\
        \midrule
        UC-01 & Đăng ký tài khoản và doanh nghiệp & Khách \\
        UC-02 & Đăng nhập & Mọi vai trò \\
        UC-03 & Hỏi đáp pháp luật & Chủ DN, Kế toán, Nhân sự \\
        UC-04 & Xem lại lịch sử hội thoại & Chủ DN, Kế toán, Nhân sự \\
        UC-05 & Tra cứu văn bản theo cây Chương--Điều & Chủ DN, Kế toán, Nhân sự \\
        UC-06 & Tìm kiếm toàn văn & Chủ DN, Kế toán, Nhân sự \\
        UC-07 & Tải lên và soát xét hợp đồng & Chủ DN, Kế toán, Nhân sự \\
        UC-08 & Xem kết quả soát xét & Chủ DN, Kế toán, Nhân sự \\
        UC-09 & Xem và cập nhật lịch tuân thủ & Chủ DN, Kế toán, Nhân sự \\
        UC-10 & Nhập bổ sung văn bản, dựng lại chỉ mục & Quản trị viên \\
        UC-11 & Quản lý người dùng & Quản trị viên \\
        UC-12 & Chạy đánh giá chất lượng truy hồi & Quản trị viên \\
        \bottomrule
    \end{tabular}
\end{table}

\subsection{Đặc tả use case}

Hai use case cốt lõi được đặc tả chi tiết dưới đây.

\begin{table}[htp]
    \centering
    \caption{Đặc tả use case UC-03 --- Hỏi đáp pháp luật.}
    \label{tab:uc03}
    \begin{tabular}{p{3.4cm}p{10.6cm}}
        \toprule
        \textbf{Tên trường} & \textbf{Nội dung} \\
        \midrule
        Use case Id & UC-03 \\
        Tên use case & Hỏi đáp pháp luật \\
        Mô tả vắn tắt & Người dùng đặt câu hỏi pháp lý bằng tiếng Việt tự nhiên
        và nhận câu trả lời kèm trích dẫn Điều luật căn cứ. \\
        Actor chính & Chủ doanh nghiệp, Kế toán, Nhân sự \\
        Actor phụ & Dịch vụ LLM và embedding bên ngoài \\
        Tiền điều kiện & Người dùng đã đăng nhập; chỉ mục vector đã sẵn sàng. \\
        Hậu điều kiện & Câu hỏi và câu trả lời kèm trích dẫn được lưu vào hội
        thoại. \\
        Luồng hoạt động & 1. Người dùng nhập câu hỏi và gửi. \newline
        2. Hệ thống tạo hội thoại mới nếu chưa có, lưu câu hỏi. \newline
        3. Hệ thống chạy pipeline bảy bước ở mục~\ref{sec:pipeline}, phát tiến độ
        sau mỗi bước. \newline
        4. Giao diện hiển thị câu trả lời tăng dần theo từng token. \newline
        5. Hệ thống đối chiếu trích dẫn với tập Điều đã truy hồi. \newline
        6. Hệ thống lưu câu trả lời kèm trích dẫn và vết xử lý. \\
        Luồng thay thế & 3a. Câu hỏi ngoài miền pháp luật: hệ thống trả lời ngay
        ở bước phân tích ý định, bỏ qua truy hồi. \newline
        3b. Không tìm được Điều liên quan: hệ thống thông báo không tìm thấy căn
        cứ thay vì trả lời chung chung. \\
        Luồng ngoại lệ & E1. Dịch vụ LLM lỗi hoặc hết hạn mức: hệ thống trả về
        thông báo lỗi rõ ràng, câu hỏi đã lưu không bị mất. \newline
        E2. Chỉ mục vector chưa sẵn sàng: hệ thống trả mã 503 kèm giải thích. \newline
        E3. Vượt hạn mức tần suất gọi: hệ thống trả mã 429. \\
        \bottomrule
    \end{tabular}
\end{table}

\begin{table}[htp]
    \centering
    \caption{Đặc tả use case UC-07 --- Tải lên và soát xét hợp đồng.}
    \label{tab:uc07}
    \begin{tabular}{p{3.4cm}p{10.6cm}}
        \toprule
        \textbf{Tên trường} & \textbf{Nội dung} \\
        \midrule
        Use case Id & UC-07 \\
        Tên use case & Tải lên và soát xét hợp đồng \\
        Mô tả vắn tắt & Người dùng tải lên hợp đồng dạng PDF hoặc DOCX; hệ thống
        phân tích và trả về danh sách rủi ro kèm căn cứ pháp lý. \\
        Actor chính & Chủ doanh nghiệp, Kế toán, Nhân sự \\
        Actor phụ & Dịch vụ LLM \\
        Tiền điều kiện & Người dùng đã đăng nhập và thuộc một doanh nghiệp. \\
        Hậu điều kiện & Tài liệu ở trạng thái \texttt{done} kèm kết quả soát xét,
        hoặc \texttt{failed} kèm lý do. \\
        Luồng hoạt động & 1. Người dùng chọn tệp và tải lên. \newline
        2. Hệ thống kiểm tra định dạng và kích thước, trích văn bản, lưu tài
        liệu ở trạng thái \texttt{pending}. \newline
        3. Người dùng yêu cầu soát xét. \newline
        4. Hệ thống chuyển trạng thái \texttt{processing} và xử lý nền: tách hợp
        đồng theo điều khoản, truy hồi quy định liên quan cho từng điều khoản,
        chấm mức rủi ro và đề xuất sửa. \newline
        5. Hệ thống tổng hợp điểm rủi ro chung, chuyển trạng thái \texttt{done}. \newline
        6. Giao diện hỏi trạng thái định kỳ và hiển thị kết quả khi hoàn tất. \\
        Luồng thay thế & 2a. Tệp là văn bản thuần: bỏ qua bước trích xuất. \newline
        4a. Hợp đồng có nhiều điều khoản hơn ngưỡng cấu hình: chỉ soát xét phần
        đầu và ghi rõ giới hạn này trong kết quả. \\
        Luồng ngoại lệ & E1. Tệp vượt dung lượng cho phép hoặc sai định dạng: từ
        chối ngay lúc tải lên. \newline
        E2. Không trích được văn bản (PDF ảnh scan): trạng thái \texttt{failed}
        kèm lý do. \newline
        E3. Lỗi trong lúc xử lý nền: trạng thái \texttt{failed}, người dùng có
        thể yêu cầu soát xét lại. \\
        \bottomrule
    \end{tabular}
\end{table}

\section{Mô hình hoá hệ thống}

\subsection{Sơ đồ lớp}

Các lớp thực thể của hệ thống được thể hiện ở Hình~\ref{fig:so-do-lop}. Chúng chia
thành ba nhóm theo nghiệp vụ: nhóm người dùng và tổ chức, nhóm hội thoại, và nhóm
tài liệu kèm tuân thủ. Nhóm thứ tư là dữ liệu pháp luật, đứng độc lập vì không
thuộc về doanh nghiệp nào.

\begin{figure}[htp]
    \centering
    \includegraphics[width=0.95\linewidth]{so-do-lop.png}
    \caption{Sơ đồ lớp thực thể.}
    \label{fig:so-do-lop}
\end{figure}

\subsection{Sơ đồ tuần tự}

Hình~\ref{fig:tuan-tu-chat} mô tả trình tự tương tác khi thực hiện UC-03. Điểm
đáng chú ý là câu hỏi của người dùng được lưu vào CSDL \textit{trước} khi gọi
pipeline, không phải sau. Thứ tự này đảm bảo câu hỏi không bị mất nếu pipeline lỗi
giữa đường, và người dùng có thể thử lại trên cùng hội thoại.

\begin{figure}[htp]
    \centering
    \includegraphics[width=0.95\linewidth]{so-do-tuan-tu-chat.png}
    \caption{Sơ đồ tuần tự use case Hỏi đáp pháp luật.}
    \label{fig:tuan-tu-chat}
\end{figure}

Hình~\ref{fig:tuan-tu-hopdong} mô tả trình tự của UC-07. Khác biệt căn bản so với
luồng hỏi đáp: yêu cầu soát xét trả về ngay sau khi ghi nhận, việc xử lý diễn ra
nền, và giao diện hỏi trạng thái định kỳ. Thiết kế này là bắt buộc vì soát xét một
hợp đồng cần chạy pipeline truy hồi cho từng điều khoản, tổng thời gian vượt xa
giới hạn hợp lý của một yêu cầu HTTP.

\begin{figure}[htp]
    \centering
    \includegraphics[width=0.95\linewidth]{so-do-tuan-tu-hop-dong.png}
    \caption{Sơ đồ tuần tự use case Soát xét hợp đồng.}
    \label{fig:tuan-tu-hopdong}
\end{figure}

\subsection{Sơ đồ hoạt động}

Hình~\ref{fig:hoat-dong-rag} thể hiện luồng xử lý của pipeline hỏi đáp cùng các
điểm rẽ nhánh: rẽ nhánh sớm khi câu hỏi ngoài miền pháp luật, và rẽ nhánh khi
không truy hồi được Điều nào liên quan.

\begin{figure}[htp]
    \centering
    \includegraphics[width=0.7\linewidth]{so-do-hoat-dong-rag.png}
    \caption{Sơ đồ hoạt động của pipeline hỏi đáp.}
    \label{fig:hoat-dong-rag}
\end{figure}

\subsection{Sơ đồ trạng thái}

Tài liệu hợp đồng là đối tượng có vòng đời rõ rệt nhất trong hệ thống, gồm bốn
trạng thái với các chuyển tiếp ở Hình~\ref{fig:trang-thai-tai-lieu}:
\texttt{pending} sau khi tải lên thành công, \texttt{processing} khi tác vụ nền
bắt đầu, rồi \texttt{done} hoặc \texttt{failed}. Từ \texttt{failed} có thể quay
lại \texttt{processing} khi người dùng yêu cầu soát xét lại --- chuyển tiếp này
cần thiết vì phần lớn nguyên nhân thất bại là lỗi tạm thời của dịch vụ bên ngoài.

\begin{figure}[htp]
    \centering
    \includegraphics[width=0.7\linewidth]{so-do-trang-thai-tai-lieu.png}
    \caption{Sơ đồ trạng thái của tài liệu hợp đồng.}
    \label{fig:trang-thai-tai-lieu}
\end{figure}

\section{Thiết kế cơ sở dữ liệu hệ thống}

\subsection{Lược đồ cơ sở dữ liệu quan hệ}

CSDL gồm mười hai bảng, liệt kê ở Bảng~\ref{tab:bang-du-lieu} và thể hiện quan hệ
ở Hình~\ref{fig:luoc-do-csdl}.

\begin{table}[htp]
    \centering
    \caption{Danh sách các bảng dữ liệu chính.}
    \label{tab:bang-du-lieu}
    \begin{tabular}{p{3.9cm}p{9.9cm}}
        \toprule
        \textbf{Bảng} & \textbf{Nội dung và các thuộc tính chính} \\
        \midrule
        \texttt{organizations} & Doanh nghiệp: tên, mã số thuế, loại hình, số lao
        động. Hai thuộc tính sau là căn cứ sinh lịch tuân thủ. \\
        \texttt{users} & Người dùng: email (duy nhất), mật khẩu đã băm, vai trò,
        doanh nghiệp, trạng thái kích hoạt. \\
        \texttt{conversations} & Hội thoại: chủ sở hữu, tiêu đề tự sinh, cờ lưu
        trữ. \\
        \texttt{messages} & Tin nhắn: hội thoại, vai trò, nội dung, trích dẫn và
        vết xử lý dạng JSONB, số thứ tự tăng đơn điệu. \\
        \texttt{documents} & Tài liệu tải lên: doanh nghiệp, người tải, tên tệp,
        đường dẫn lưu trữ, văn bản đã trích, trạng thái. \\
        \texttt{contract\_reviews} & Lượt soát xét: tài liệu, trạng thái, điểm rủi
        ro tổng hợp, bản tóm tắt. \\
        \texttt{contract\_findings} & Phát hiện rủi ro: điều khoản, mức rủi ro,
        giải thích, đề xuất sửa, Điều luật căn cứ. \\
        \texttt{compliance\_rules} & Nghĩa vụ định kỳ: mã, tiêu đề, tần suất, quy
        tắc tính hạn, căn cứ pháp lý, điều kiện áp dụng. \\
        \texttt{compliance\_tasks} & Tác vụ sinh ra cho từng doanh nghiệp: nghĩa
        vụ, hạn chót, trạng thái, thời điểm hoàn thành. \\
        \texttt{legal\_knowledge\_records} & Corpus: mỗi dòng một Điều luật kèm
        metadata và cột \texttt{tsvector} phục vụ tìm kiếm toàn văn. \\
        \texttt{laws} & Danh mục văn bản, dẫn xuất từ corpus, phục vụ hiển thị
        danh sách và bộ lọc. \\
        \texttt{audit\_logs} & Vết thao tác: người dùng, hành động, đối tượng. \\
        \bottomrule
    \end{tabular}
\end{table}

\begin{figure}[htp]
    \centering
    \includegraphics[width=0.95\linewidth]{luoc-do-csdl.png}
    \caption{Lược đồ cơ sở dữ liệu quan hệ.}
    \label{fig:luoc-do-csdl}
\end{figure}

\subsection{Các mối quan hệ quan trọng}

Bốn quan hệ quyết định nghiệp vụ cốt lõi:

\textbf{organizations 1--n users.} Một doanh nghiệp có nhiều người dùng thuộc các
vai trò khác nhau. Đây là quan hệ nền tảng cho cách ly dữ liệu: mọi truy vấn về
hợp đồng và lịch tuân thủ đều lọc theo doanh nghiệp của người dùng đang đăng nhập.

\textbf{users 1--n conversations 1--n messages.} Hội thoại thuộc về từng người
dùng chứ không thuộc về doanh nghiệp --- lựa chọn có chủ ý, vì câu hỏi pháp lý của
một cá nhân có thể mang tính riêng tư và không nên hiển thị cho đồng nghiệp.

\textbf{documents 1--n contract\_reviews 1--n contract\_findings.} Một tài liệu có
thể được soát xét nhiều lần, mỗi lượt sinh ra nhiều phát hiện. Giữ nhiều lượt soát
xét thay vì ghi đè cho phép so sánh kết quả trước và sau khi sửa hợp đồng.

\textbf{compliance\_rules n--n organizations qua compliance\_tasks.} Nghĩa vụ được
định nghĩa một lần dùng chung cho mọi doanh nghiệp; tác vụ cụ thể mới gắn với từng
doanh nghiệp kèm hạn chót riêng. Tách như vậy để khi bổ sung hay sửa một nghĩa vụ
thì chỉ cần thay đổi ở một chỗ.

\subsection{Các ràng buộc dữ liệu}

\textbf{Khóa chính và khóa ngoại.} Các bảng nghiệp vụ dùng khóa chính là số nguyên
tự tăng; riêng \texttt{laws} dùng số hiệu văn bản làm khóa chính vì đây vốn là
định danh tự nhiên và duy nhất. Khóa ngoại đặt hành vi xóa lan truyền cho quan hệ
sở hữu thực sự: xóa hội thoại thì xóa tin nhắn của nó, xóa tài liệu thì xóa các
lượt soát xét --- những dữ liệu này không có ý nghĩa khi tách khỏi bản ghi cha.

\textbf{Ràng buộc duy nhất.} Email người dùng là duy nhất toàn hệ thống. Mã nghĩa
vụ tuân thủ là duy nhất. Cặp (doanh nghiệp, nghĩa vụ, hạn chót) là duy nhất, nhằm
ngăn sinh trùng tác vụ khi người dùng bấm tạo lại lịch nhiều lần.

\textbf{Ràng buộc miền giá trị.} Vai trò người dùng, trạng thái tài liệu, mức rủi
ro và tần suất nghĩa vụ đều dùng kiểu liệt kê thay vì chuỗi tự do, để CSDL từ chối
giá trị sai ngay ở tầng lưu trữ.

\textbf{Thứ tự tin nhắn.} Bảng \texttt{messages} có cột số thứ tự tăng đơn điệu do
CSDL sinh, và việc sắp xếp tin nhắn dựa trên cột này chứ không dựa trên thời điểm
tạo. Nguyên nhân là hàm lấy thời gian hiện tại của PostgreSQL trả về thời điểm bắt
đầu giao dịch, trong khi câu hỏi và câu trả lời được ghi trong cùng một giao dịch
--- nếu sắp theo thời điểm tạo thì hai tin nhắn có cùng giá trị và thứ tự hiển thị
trở nên không xác định.

\section{Thiết kế tìm kiếm toàn văn tiếng Việt}

Chức năng tìm kiếm toàn văn có hai vấn đề riêng của tiếng Việt, đáng trình bày
tách thành một mục vì cách giải quyết không hiển nhiên.

\textbf{Vấn đề thứ nhất: PostgreSQL không có từ điển tiếng Việt.} Các cấu hình
tìm kiếm sẵn có đều gắn với một ngôn ngữ cụ thể và không có cấu hình cho tiếng
Việt. Giải pháp của hệ thống là bỏ dấu toàn bộ văn bản rồi dùng cấu hình
\texttt{simple}. Cách này đồng thời xử lý luôn thói quen gõ không dấu của người
dùng: truy vấn \textit{``hop dong lao dong''} và văn bản \textit{``hợp đồng lao
động''} sau khi bỏ dấu trở thành cùng một chuỗi.

Việc hiện thực gặp một trở ngại kỹ thuật: PostgreSQL chỉ cho phép dùng hàm ở mức
\texttt{IMMUTABLE} trong cột sinh tự động và chỉ mục, nhưng hàm bỏ dấu tiêu chuẩn
lại ở mức \texttt{STABLE} do phụ thuộc vào từ điển có thể thay đổi. Hệ thống giải
quyết bằng cách bọc hàm này trong một hàm tự định nghĩa được khai báo
\texttt{IMMUTABLE}, từ đó tạo được cột \texttt{tsvector} sinh tự động cùng chỉ mục
GIN trên nó.

\textbf{Vấn đề thứ hai: bôi đậm từ khóa mà không mất dấu.} Cách thông thường là
dùng hàm tạo đoạn trích của PostgreSQL, nhưng ở đây hàm đó chạy trên cột đã bỏ
dấu, nên đoạn trích trả về là văn bản không dấu --- không thể hiển thị cho người
đọc. Hệ thống chuyển việc bôi đậm sang phía client: backend chỉ trả về danh sách
từ khóa đã tách, client tự dò và bôi đậm trên văn bản gốc còn dấu, so khớp theo
kiểu bỏ qua dấu. Kết quả là người dùng gõ không dấu vẫn thấy đúng đoạn có dấu được
bôi đậm.

Hệ thống áp dụng ba chiến lược tìm kiếm theo thứ tự ưu tiên: tra trực tiếp khi
truy vấn chứa số Điều, xếp hạng toàn văn cho trường hợp thông thường, và so khớp
gần đúng theo n-gram để cứu các truy vấn gõ sai chính tả.

\section{Thiết kế bảo mật và phân quyền}

\textbf{Xác thực.} Mật khẩu được băm bằng argon2. Hệ thống phát hai loại token:
access token thời hạn ngắn dùng cho mọi yêu cầu, và refresh token thời hạn dài chỉ
dùng để lấy access token mới. Hai loại được phân biệt bằng một trường trong nội
dung token, nhờ đó refresh token không thể dùng thay access token để truy cập API
--- nếu không có trường phân biệt này, một token thời hạn dài bị lộ sẽ mở toàn bộ
hệ thống.

\textbf{Phân quyền.} Mỗi endpoint khai báo vai trò được phép truy cập. Các endpoint
quản trị chỉ mở cho vai trò quản trị viên.

\textbf{Cách ly dữ liệu.} Truy vấn hợp đồng và lịch tuân thủ luôn lọc theo doanh
nghiệp; truy vấn hội thoại luôn lọc theo người dùng. Khi truy cập tài nguyên của
người khác, hệ thống trả về mã 404 chứ không phải 403 --- vì mã 403 gián tiếp xác
nhận tài nguyên đó tồn tại.

\textbf{Giới hạn tần suất.} Các endpoint hỏi đáp và tải tệp bị giới hạn số lần gọi
theo đơn vị thời gian, nhằm tránh một người dùng làm cạn hạn mức của dịch vụ LLM.

\section{Thiết kế giao diện hệ thống}

Giao diện là một SPA tối giản, ưu tiên mật độ thông tin pháp lý hơn hiệu ứng.
Bốn nghiệp vụ chính dùng chung khung điều hướng bên trái; trang quản trị chỉ hiện
với vai trò quản trị viên.

\subsection{Hỏi đáp pháp luật}

Trang chat gồm ba vùng: thanh lịch sử hội thoại, luồng tin nhắn, và ô nhập.
Khi người dùng gửi câu hỏi, giao diện hiện ngay lượt hỏi rồi nhận ba loại sự kiện
SSE: tiến độ từng bước của pipeline (hiện thành vết xử lý có thể thu gọn), token
câu trả lời, và kết quả cuối kèm danh sách Điều luật. Mỗi Điều là một liên kết tới
trang tra cứu. Nút ``Hỏi AI về Điều này'' ở trang văn bản bơm sẵn câu hỏi sang đây
kèm ngữ cảnh số hiệu văn bản.

\begin{figure}[htp]
    \centering
    \includegraphics[width=0.85\linewidth]{man-hinh-chat.png}
    \caption{Giao diện hỏi đáp pháp luật: tin nhắn, vết xử lý và trích dẫn.}
    \label{fig:man-hinh-chat}
\end{figure}

\subsection{Tra cứu văn bản}

Trang danh mục lọc theo lĩnh vực và loại văn bản; trang chi tiết dựng cây
Chương~$\rightarrow$~Điều. Ô tìm kiếm toàn văn bôi đậm từ khóa trên văn bản gốc
còn dấu, kể cả khi người dùng gõ không dấu --- hệ quả của chiến lược ở mục trên.

\begin{figure}[htp]
    \centering
    \includegraphics[width=0.85\linewidth]{man-hinh-tra-cuu.png}
    \caption{Giao diện tra cứu văn bản theo cây Chương--Điều.}
    \label{fig:man-hinh-tra-cuu}
\end{figure}

\subsection{Soát xét hợp đồng và lịch tuân thủ}

Trang hợp đồng cho phép kéo-thả PDF/DOCX, theo dõi trạng thái
\texttt{pending~$\rightarrow$~processing~$\rightarrow$~done/failed}, rồi mở bảng
phát hiện rủi ro. Mỗi phát hiện gắn mức \textit{cao / trung bình / thấp}, đề xuất
sửa và Điều luật căn cứ.

Trang tuân thủ gồm bảng thống kê (sắp hết hạn, quá hạn, đã xong) và danh sách tác
vụ theo hạn. Người dùng đánh dấu hoàn thành tại chỗ; căn cứ pháp lý của mỗi nghĩa
vụ là liên kết sang trang tra cứu.

\section{Kết quả triển khai hệ thống}

\subsection{Triển khai hệ thống}

Hệ thống được đóng gói bằng Docker Compose gồm ba dịch vụ trên mạng cầu:
PostgreSQL~17, backend FastAPI và frontend nginx phục vụ SPA đồng thời chuyển tiếp
\texttt{/api}. Mỗi dịch vụ có healthcheck; backend chỉ nhận yêu cầu sau khi
PostgreSQL sẵn sàng. Dữ liệu bền vững nằm trên ba volume: CSDL, chỉ mục vector và
tệp hợp đồng tải lên.

Môi trường phát triển chạy backend bằng Python~3.10 (ràng buộc của thư viện tách
từ tiếng Việt) và frontend bằng Vite, proxy \texttt{/api} sang cổng backend.

\subsection{Chức năng đã hoàn thành}

Toàn bộ mục tiêu cụ thể ở Chương~1 đã được hiện thực: hỏi đáp có trích dẫn và
streaming, tra cứu cây Chương--Điều kèm full-text tiếng Việt, soát xét hợp đồng
nền, lịch tuân thủ sinh theo hồ sơ doanh nghiệp, xác thực JWT/argon2 với bốn vai
trò, trang quản trị corpus và lab thi đấu, đóng gói Docker Compose, kèm
\soTest{} test tự động.

\subsection{Hình ảnh chạy thực tế hệ thống}

Ba dịch vụ khởi động bằng một lệnh \texttt{docker compose up}. Khi chưa có API
key Gemini, đăng nhập, tra cứu văn bản và lịch tuân thủ vẫn hoạt động; hỏi đáp
trả mã 503 kèm hướng dẫn. Sau khi bổ sung key, gọi reindex là chức năng hỏi đáp
bật lại, không cần dựng lại container.

\subsection{Hạn chế còn tồn tại}

Hệ thống phụ thuộc hạn mức API của nhà cung cấp LLM: corpus \soDieu{} Điều không
thể nhúng lại thường xuyên trên tầng miễn phí. Soát xét hợp đồng chưa xử lý PDF
ảnh scan (không có tầng OCR). Lịch tuân thủ chưa theo dõi thay đổi pháp luật theo
thời gian --- văn bản mới phải nhập tay qua trang quản trị. Giao diện tối ưu cho
máy tính, chưa có ứng dụng di động riêng.


% ============================================================
% CHƯƠNG 4. KẾT LUẬN VÀ HƯỚNG PHÁT TRIỂN
% ============================================================
\chapter{Kết luận và hướng phát triển}

\section{Kết luận}

Đồ án đã xây dựng được một ứng dụng web trợ lý pháp lý cho DNNVV, đáp ứng bốn
nghiệp vụ đặt ra: hỏi đáp có trích dẫn tới từng Điều, tra cứu văn bản theo cấu
trúc và theo từ khóa, soát xét rủi ro hợp đồng, và nhắc nghĩa vụ tuân thủ định kỳ.
Điểm khác biệt so với chatbot đa dụng là mọi thông tin pháp lý đều truy vết được
về văn bản trong corpus; khi không tìm thấy căn cứ, hệ thống nói rõ thay vì trả
lời chung chung.

Về kỹ thuật, đồ án cho thấy hybrid retrieval (vector + BM25 + RRF) kết hợp tinh
lọc bằng LLM là cách phù hợp với khoảng cách từ vựng giữa câu hỏi thường ngày và
văn bản luật; đồng thời chỉ ra rằng full-text tiếng Việt trên PostgreSQL phải tự
thiết kế (bỏ dấu + \texttt{simple} + trigram) chứ không dùng được từ điển có sẵn.
Lịch sử hội thoại đặt ở PostgreSQL thay vì bộ nhớ tiến trình, nên hệ thống chịu
được khởi động lại và chạy nhiều worker.

Sản phẩm ứng dụng dùng Gemini qua lớp tương thích OpenAI; kết quả thi đấu gốc dùng
mô hình mã nguồn mở dưới 14B tham số. Hai con số đánh giá không so sánh trực tiếp
được và đã được tách rõ trong tài liệu.

\section{Hướng phát triển}

Một số hướng mở rộng có tính khả thi sau đồ án:

\begin{itemize}
    \item Bổ sung OCR cho PDF ảnh scan trước khi soát xét hợp đồng.
    \item Theo dõi hiệu lực văn bản theo thời gian (văn bản sửa đổi, hết hiệu lực)
    để corpus không bị đóng băng tại thời điểm thu thập.
    \item Thử nghiệm đồ thị tri thức pháp lý (GraphRAG) để khai thác viện dẫn chéo
    đã có sẵn trong metadata từng Điều.
    \item Đưa soát xét hợp đồng sang hàng đợi tác vụ phân tán khi quy mô vượt quá
    \texttt{BackgroundTasks}.
    \item Đánh giá lại F2 trên bộ test nội bộ sau khi đổi embedding, để chương
    thực nghiệm có số liệu so sánh giữa Qwen và Gemini.
\end{itemize}


% ============================================================
% TÀI LIỆU THAM KHẢO
% ============================================================
\clearpage
\phantomsection
\addcontentsline{toc}{chapter}{TÀI LIỆU THAM KHẢO}
\begin{thebibliography}{9}

\bibitem{thanh2024}
Dương Hữu Thành, \textit{Công nghệ phần mềm}, NXB Thông tin \& Truyền thông, 2024.

\bibitem{duong2026}
Duong, H.T., Bui, T.P., Nguyen, D.H., ``QLoRA and RAG Application in Building Admissions Support Chatbot: Performance Evaluation and Knowledge Update Ability,'' In: Dao, N.N., Truong Hoang, V., Dornaika, F. (eds) \textit{Explainable Intelligence in Digital Twins}. EIDT 2025. Lecture Notes in Electrical Engineering, vol 1531. Springer, Singapore, 2026.

\bibitem{lewis2020rag}
Lewis, P., Perez, E., Piktus, A., et al., ``Retrieval-Augmented Generation for
Knowledge-Intensive NLP Tasks,'' \textit{Advances in Neural Information
Processing Systems (NeurIPS)}, vol. 33, pp. 9459--9474, 2020.

\bibitem{karpukhin2020dpr}
Karpukhin, V., Oğuz, B., Min, S., et al., ``Dense Passage Retrieval for
Open-Domain Question Answering,'' \textit{Proceedings of EMNLP 2020},
pp. 6769--6781, 2020.

\bibitem{robertson2009bm25}
Robertson, S., Zaragoza, H., ``The Probabilistic Relevance Framework: BM25 and
Beyond,'' \textit{Foundations and Trends in Information Retrieval}, vol. 3,
no. 4, pp. 333--389, 2009.

\bibitem{cormack2009rrf}
Cormack, G.V., Clarke, C.L.A., Buettcher, S., ``Reciprocal Rank Fusion
Outperforms Condorcet and Individual Rank Learning Methods,''
\textit{Proceedings of SIGIR 2009}, pp. 758--759, 2009.

\bibitem{gao2023hyde}
Gao, L., Ma, X., Lin, J., Callan, J., ``Precise Zero-Shot Dense Retrieval without
Relevance Labels,'' \textit{Proceedings of ACL 2023}, pp. 1762--1777, 2023.

\bibitem{quochoi2017}
Quốc hội, \textit{Luật Hỗ trợ doanh nghiệp nhỏ và vừa}, số 04/2017/QH14, 2017.

\bibitem{postgresql2024}
PostgreSQL Global Development Group, ``PostgreSQL Documentation: Full Text
Search.'' [Online]. Available:
\url{https://www.postgresql.org/docs/current/textsearch.html}.

\bibitem{openai2026}
OpenAI, ``ChatGPT.'' [Online]. Available: \url{https://chat.openai.com}. [Accessed: Jul. 1, 2026].

\end{thebibliography}


% ============================================================
% PHỤ LỤC
% ============================================================
\clearpage
\phantomsection
\addcontentsline{toc}{chapter}{PHỤ LỤC}
\chapter*{PHỤ LỤC}

\textbf{A. Chạy hệ thống.} Sao chép \texttt{.env.example} thành \texttt{.env},
điền \texttt{LLM\_API\_KEY} và \texttt{JWT\_SECRET\_KEY}, rồi
\texttt{docker compose up -d --build}. Lần đầu nạp corpus:
\texttt{docker compose exec backend python scripts/load\_postgres.py --truncate}
và khởi động lại backend.

\textbf{B. Kiểm thử.} Backend: \texttt{cd backend \&\& uv run pytest}
(PostgreSQL thật, database \texttt{legal\_assistant\_test}). Frontend:
\texttt{cd frontend \&\& npm test}.

\textbf{C. Lab thi đấu.} \texttt{POST /api/v1/lab/competition} (vai trò admin)
giữ nguyên format \texttt{results.json} để tái tạo số liệu đánh giá.

\end{document}